% ==============================================================================
% CHƯƠNG 2 - PHẦN D: CHUYỂN ĐỔI, PHÂN MẢNH VÀ CHỈ MỤC (TRANSFORM & INDEX)
% ==============================================================================

\section{Kiến Trúc Lưu Trữ Cột Zero-Local-Disk và Compaction LSM-Style}
\label{sec:ch2_zstd_compaction}

\subsection{Ghi Bất Biến Theo Lô Nhỏ (\texorpdfstring{$O(\text{batch})$}{O(batch)}) trên Bộ Đệm RAM}
\label{sec:ch2_append_part}

Toàn bộ tầng lưu trữ trung gian sử dụng lớp \texttt{LakehouseIO} (\path{src/storage/lakehouse_io.py}) để tuần tự hóa bảng dữ liệu thành tệp Apache Parquet nén \textbf{ZSTD} trực tiếp qua bộ nhớ đệm \texttt{io.BytesIO()} trước khi tải thẳng lên MinIO S3 (\textit{Zero Local Disk footprint}). Khi worker \texttt{fetch} hoặc \texttt{embed} xử lý mỗi micro-batch (\texttt{16} hoặc \texttt{32} bài), việc đọc-sửa-ghi đè toàn bộ tệp lớn theo năm (\path{papers.parquet} hoặc \path{chunks.parquet}) sẽ gây độ phức tạp $O(N)$ và nghẽn mạng. Thay vào đó, cả \texttt{SilverLakehouseWriter.append\_part()} (\path{src/transformation/silver_writer.py}) và \texttt{ChunkStore.append\_part()} (\path{src/storage/chunk_store.py}) ghi các phân đoạn nhỏ \textbf{bất biến} với độ phức tạp $O(\text{batch})$:
\begin{itemize}
  \item Bảng tài liệu: \path{silver/papers/year={YYYY}/parts/part-{YYYYMMDDTHHMMSS}-{uuid8}.parquet}
  \item Bảng phân đoạn chunk: \path{silver/chunks/year={YYYY}/parts/part-{YYYYMMDDTHHMMSS}-{uuid8}.parquet}
\end{itemize}

\subsection{Cơ Chế Compaction Chu Kỳ 3 Giờ (\texttt{dag\_compact\_silver}) và Khóa Đồng Thời}
\label{sec:ch2_compaction}

Một đợt thu thập Tier 2 có thể sinh ra xấp xỉ \texttt{2,000} tệp \texttt{part-*} mỗi ngày; nếu để phân mảnh quá lâu, các truy vấn DuckDB S3 trên tầng Silver phải mở hàng nghìn kết nối HTTP khiến độ trễ tăng cao. DAG \path{dags/dag_compact_silver.py} kích hoạt định kỳ mỗi 3 giờ (\texttt{0 */3 * * *}) gọi \texttt{POST /api/ingestion/compact-silver} (\path{backend/app/api/endpoints/ingestion.py}), thực thi theo cơ chế kiểu \textbf{LSM-Tree Compaction} an toàn bên cạnh worker đang chạy:
\begin{enumerate}[label=\textbf{\arabic*.}]
  \item \textbf{Chốt chặn đơn tiến trình} (\path{_compact_lock}): Endpoint dùng \path{threading.Lock().acquire(blocking=False)} để trả ngay mã \texttt{HTTP 409 Conflict} nếu một lượt compaction khác đang thực thi.
  \item \textbf{Hợp nhất \path{silver/papers/} (\texttt{SilverLakehouseWriter.compact\_all})}: Liệt kê danh sách khóa \texttt{part\_keys} hiện hữu tại từng năm, hợp nhất cùng \path{papers.parquet}, khử trùng theo khóa \texttt{paper\_id} (\texttt{keep="last"}), ghi tệp mới lên S3 và \textbf{chỉ gọi \texttt{delete\_objects\_batch(part\_keys)} sau khi tệp gộp đã ghi thành công}. Các tệp \texttt{part-*} mới sinh ra trong lúc compaction đang chạy không nằm trong danh sách chụp ban đầu nên được giữ nguyên cho chu kỳ 3 giờ kế tiếp.
  \item \textbf{Hợp nhất \path{silver/chunks/} (\texttt{ChunkStore.compact\_all} và \texttt{latest\_only})}: Đọc bảng cũ cùng các \texttt{part\_keys}, chạy hàm \texttt{ChunkStore.latest\_only()} (\path{src/storage/chunk_store.py}) để giữ lại đúng phiên bản mới nhất (\texttt{max(indexed\_at)}) cho từng \texttt{paper\_id} (xóa sạch dấu vết chunk cũ của các bài đã \texttt{refresh=True}), chọn dòng cuối cùng cho mỗi \texttt{chunk\_id}, rồi \textbf{sắp xếp tăng dần toàn bộ bảng theo \texttt{chunk\_id}} và ghi xuất bằng nén ZSTD với kích thước nhóm dòng nhỏ \texttt{ROW\_GROUP\_SIZE = 512} ($\sim 0.3\text{ MB}$ sau nén mỗi group).
\end{enumerate}

\section{Thiết Kế \texttt{ChunkStore} và Thủy Hóa Byte-Range Trên S3 (\texttt{D-058})}
\label{sec:ch2_chunkstore}

\subsection{Kiến Trúc Tách Payload Qdrant (\texttt{D-058}) và Lược Đồ 14 Cột}
\label{sec:ch2_d058}

Trong giai đoạn đầu, toàn bộ văn bản phân đoạn (\texttt{text}, \texttt{context\_text}), danh sách tác giả và tiêu đề bài báo được nhồi thẳng vào payload JSON của từng điểm vector trong Qdrant. Khi quy mô vượt mốc \texttt{1,000,000} vector, thiết kế này gây phình bộ nhớ RAM của Qdrant, nhân đôi dung lượng lưu trữ với bảng Silver và làm chậm thao tác quét điểm (\textit{scroll}). Theo quyết định kiến trúc \textbf{\texttt{D-058}} (trích đoạn triển khai tại Đoạn mã~\ref{lst:payload_hydrate} trong Phần Phụ lục), toàn bộ văn bản và siêu dữ liệu đầy đủ được tách về bảng cột \path{silver/chunks/} (\texttt{ChunkStore}), còn payload trong Qdrant chỉ giữ \textbf{5 trường tối thiểu} (\texttt{QDRANT\_PAYLOAD\_FIELDS = ("chunk\_id", "paper\_id", "year", "primary\_category", "section\_type")}). Bảng~\ref{tab:ch2_d058_compare} đối chiếu hai phương án thiết kế và Bảng~\ref{tab:chunkstore_schema} đặc tả chi tiết 14 cột của \texttt{ChunkStore.SCHEMA}.

\begin{table}[htbp]
  \centering
  \small
  \caption[So sánh kiến trúc lưu trữ vector trước và sau quyết định D-058]{So sánh kiến trúc lưu trữ vector trước và sau quyết định kiến trúc \texttt{D-058}}
  \label{tab:ch2_d058_compare}
  \setlength{\tabcolsep}{3.5pt}
  \begin{tabularx}{\linewidth}{|>{\raggedright\arraybackslash}p{31mm}|>{\raggedright\arraybackslash}p{51mm}|L|}
    \hline
    \textbf{Chiều so sánh} & \textbf{Kiến trúc cũ (Full Payload trong Vector DB)} & \textbf{Kiến trúc \texttt{D-058} (\texttt{Qdrant Minimal} + \texttt{ChunkStore})} \\
    \hline
    \textbf{Payload trong Qdrant} & Vector 768-D + toàn bộ \path{text}, \path{context_text}, \path{title}, \path{authors}, \path{doi} & Vector 768-D + \textbf{5 trường lọc}: \path{chunk_id}, \path{paper_id}, \path{year}, \path{primary_category}, \path{section_type} \\
    \hline
    \textbf{Nơi giữ văn bản gốc} & Nhân bản song song trong Qdrant / LanceDB & Nguồn chân lý duy nhất tại \path{silver/chunks/year=YYYY/chunks.parquet} (nén ZSTD) \\
    \hline
    \textbf{Đọc kết quả Top-$k$} & Trả trực tiếp từ payload vector & Qdrant trả Top-$k$ ID $\rightarrow$ \path{ChunkStore.hydrate()} đọc dải byte S3 có LRU cache \\
    \hline
    \textbf{Rủi ro \& Kiểm soát} & Tốn RAM chỉ mục vector khi đạt $>1\text{M}$ điểm & Nếu quên hydrate, reranker nhận chuỗi rỗng ($0\%$ delta); khắc phục bằng ghi Silver trước Qdrant + hydrate trong \path{QdrantManager._hydrate()} và \path{RetrievalService} \\
    \hline
  \end{tabularx}
  \ReportNote{Nguồn: Đối chiếu \path{src/storage/chunk_store.py}, \path{src/indexing/qdrant_manager.py} và báo cáo \path{docs/benchmarks/laya_pagerank_benchmark_report.md} (quyết định \texttt{D-058}).}
\end{table}

\begin{table}[htbp]
  \centering
  \small
  \caption[Lược đồ 14 cột của bảng phân đoạn ChunkStore]{Lược đồ chuẩn hóa 14 cột của bảng \texttt{silver/chunks/} (\texttt{ChunkStore.SCHEMA})}
  \label{tab:chunkstore_schema}
  \setlength{\tabcolsep}{3.5pt}
  \begin{tabularx}{\linewidth}{|>{\centering\arraybackslash}p{8mm}|>{\raggedright\arraybackslash}p{28mm}|>{\raggedright\arraybackslash}p{24mm}|L|}
    \hline
    \textbf{STT} & \textbf{Tên cột} & \textbf{Kiểu Arrow} & \textbf{Ý nghĩa \& Vai trò thủy hóa / khai phá} \\
    \hline
    1 & \path{chunk_id} & \path{pa.string()} & Khóa chính phân đoạn (\texttt{\{pid\}\_abs}, \texttt{\{pid\}\_sec\{idx\}\_p\{part\}}), khóa sắp xếp \& thống kê Min/Max Row Group \\
    \hline
    2 & \path{paper_id} & \path{pa.string()} & Mã định danh bài báo gốc dùng nhóm phiên bản theo bài trong \path{latest_only()} \\
    \hline
    3 & \path{title} & \path{pa.string()} & Tiêu đề bài báo dùng ghép ngữ cảnh cho Cross-Encoder reranker và hiển thị citation \\
    \hline
    4 & \path{authors} & \path{pa.list_(pa.string())} & Danh sách tác giả truyền vào bộ khớp uy tín PageRank (\path{_match_authority}) \\
    \hline
    5 & \path{primary_category} & \path{pa.string()} & Nhãn chuyên ngành chính (đồng bộ với trường lọc Qdrant) \\
    \hline
    6 & \path{year} & \path{pa.int64()} & Năm công bố kiểu nguyên (trích qua \path{_parse_year}), hỗ trợ phân vùng và lọc thời gian \\
    \hline
    7 & \path{section_title} & \path{pa.string()} & Tên mục cụ thể (vd. \texttt{3. Methodology (Part 1)}), dùng cho gắn thẻ \texttt{[Paper: ID, Section: \ldots]} \\
    \hline
    8 & \path{section_type} & \path{pa.string()} & Nhãn chuẩn hóa 9 loại mục (\path{abstract}, \path{methodology}, \path{experiments}, \ldots) \\
    \hline
    9 & \path{text} & \path{pa.string()} & Toàn văn nội dung đoạn trích (không lặp tiền tố \texttt{Paper:} để trình bày sạch trên UI \& Prompt) \\
    \hline
    10 & \path{word_count} & \path{pa.int64()} & Số từ thực tế của đoạn trích (\path{len(words)}) \\
    \hline
    11 & \path{source} & \path{pa.string()} & Nguồn thu thập bài báo (\path{arxiv}, \path{cvf}, \path{semantic_scholar}, \ldots) \\
    \hline
    12 & \path{tier} & \path{pa.string()} & Tầng thu thập (\texttt{tier1}, \texttt{tier2}, \texttt{tier3}) \\
    \hline
    13 & \path{content_origin} & \path{pa.string()} & Nguồn gốc nội dung (\path{arxiv_html5}, \path{ar5iv_html5}, \path{open_access_pdf}, \path{abstract_fallback}) \\
    \hline
    14 & \path{indexed_at} & \path{pa.string()} & Dấu thời gian UTC ISO 8601 khi lập chỉ mục; phân giải xung đột ghi mới nhất khi \path{refresh=True} \\
    \hline
  \end{tabularx}
  \ReportNote{Nguồn: Định nghĩa \texttt{ChunkStore.SCHEMA} và hàm \texttt{to\_table()} tại \path{src/storage/chunk_store.py} (kiểm chứng trên commit \texttt{f421707}/\texttt{aae74de}, ngày 25/11/2026); dùng cho thủy hóa văn bản phân đoạn và đối khớp PageRank tách khỏi payload Qdrant theo quyết định \texttt{D-058}.}
\end{table}


\subsection{Thuật Toán Đọc Dải Byte Song Song và Bộ Nhớ Đệm LRU}
\label{sec:ch2_byte_range}

Thách thức lớn nhất của việc tách văn bản về MinIO S3 là làm sao thủy hóa \texttt{10--40} \texttt{chunk\_id} trong vài chục mili-giây mà không phải tải cả tệp Parquet hàng trăm megabyte. Lớp \texttt{ChunkStore} giải quyết bằng kiến trúc truy xuất 5 bước (Hình~\ref{fig:chunkstore_arch}):
\begin{enumerate}[label=\textbf{\arabic*.}]
  \item \textbf{Bộ nhớ đệm LRU trong RAM (\texttt{\_cache})}: Cấu trúc \texttt{collections.OrderedDict} lưu tối đa \texttt{20,000} bản ghi chunk nóng, bảo vệ truy cập đa luồng bằng \texttt{threading.Lock()} và đẩy mục vừa đọc lên cuối (\texttt{move\_to\_end()}).
  \item \textbf{Cache Footer và Cắt tỉa Row Group theo Thống kê Min/Max (\texttt{\_file\_meta})}: Khi một tệp Parquet (hoặc phiên bản \texttt{(key, size, last\_modified)} mới) được tiếp cận lần đầu, \texttt{ChunkStore} chỉ đọc phần \textbf{Parquet Footer} để trích xuất siêu dữ liệu các Row Group (\texttt{ROW\_GROUP\_SIZE = 512} dòng $\approx 0.3\text{ MB}$ nén mỗi nhóm) và lưu mảng cặp cận \texttt{(min, max)} của cột \texttt{chunk\_id}. Vì \path{chunks.parquet} sau compaction được sắp xếp tăng dần theo \texttt{chunk\_id}, tập khoảng \texttt{(min, max)} gần như không giao nhau, cho phép loại bỏ $>99\%$ Row Group ngay trên bộ nhớ mà không tốn một byte truyền tải mạng nào.
  \item \textbf{Đọc Dải Byte Song Song Gom Lô} (\path{pre_buffer=True}, \path{max_workers=16}): Đối với một số ít cặp \texttt{(file, row\_group)} thỏa điều kiện $\text{lo} \le \text{cid} \le \text{hi}$, hệ thống điều phối trên \path{ThreadPoolExecutor(max_workers=16)}. Mỗi tác vụ gọi \path{pq.ParquetFile(fh, metadata=meta, pre_buffer=True).read_row_group(group)} để PyArrow tự động gộp các khoảng byte liền kề thành một yêu cầu HTTP S3 Range Read duy nhất, sau đó lọc chính xác bằng nhân C++ \path{pc.is_in(t["chunk_id"], value_set=want)}.
  \item \textbf{Cơ chế Thử lại Danh sách Hai lượt (\textit{Double-Pass Relist})}: Danh sách tệp trên S3 được cache \texttt{list\_ttl = 30.0 s}. Nếu truy vấn gặp một \texttt{chunk\_id} mới vừa được worker \texttt{embed} ghi trong vòng vài giây trước đó (\texttt{attempt == 0} bị hụt), \texttt{get()} tự động làm mới danh sách tệp (\texttt{force=True}) ở lượt thứ hai để phát hiện các \texttt{part-*} vừa sinh.
  \item \textbf{Kênh Dự phòng Địa phương và Theo yêu cầu (\texttt{\_lookup\_local\_silver})}: Khi một điểm thuộc tập kiểm thử lịch sử vắng mặt trên kho S3 chunks hiện tại, \texttt{hydrate()} tự động tra cứu vào các bảng Silver cục bộ (\path{data/silver/...}) hoặc gọi API arXiv theo yêu cầu (\texttt{\_fetch\_arxiv\_on\_demand()}, có lưu đệm \path{data/silver/arxiv_metadata_cache.json}).
\end{enumerate}

\begin{figure}[htbp]
  \centering
  \DrawChunkStoreArchitecture
  \caption[Kiến trúc ChunkStore D-058: Ghi part ZSTD, Compaction và Đọc dải byte S3]{Kiến trúc \texttt{ChunkStore} (\texttt{D-058}): Ghi part ZSTD, Compaction LSM 3 giờ và Thủy hóa S3 Byte-Range song song}
  \ReportNote{Nguồn: Thiết kế và cài đặt tại \path{src/storage/chunk_store.py} cùng lịch compaction \path{dags/dag_compact_silver.py}.}
  \label{fig:chunkstore_arch}
\end{figure}

\subsection{Công Cụ Di Trú Payload và Bù Khuyết Chỉ Mục}
\label{sec:ch2_migration_tools}

Để chuyển đổi dữ liệu hiện hữu sang kiến trúc \texttt{D-058} mà không mất mát lịch sử, thư mục \path{scripts/} cung cấp bộ ba công cụ kiểm toán và đồng bộ an toàn:
\begin{itemize}
  \item \path{scripts/migrate_chunks_to_silver.py}: Thực thi 4 bước tuần tự \texttt{export} (quét payload Qdrant thành lô \texttt{50,000} dòng ghi vào \path{silver/chunks/}) $\rightarrow$ \texttt{compact} $\rightarrow$ \texttt{verify} (đối chiếu tập \texttt{chunk\_id} giữa Qdrant và Silver đảm bảo \texttt{missing\_in\_silver == 0} và không trùng lặp) $\rightarrow$ \texttt{--strip} (chỉ xóa các khóa payload cũ ngoài \texttt{QDRANT\_PAYLOAD\_FIELDS} sau khi bước \texttt{verify} vượt qua 100\%).
  \item \path{scripts/backfill_missing_papers.py}: Phát hiện các bài trong \path{silver/papers/} chưa có điểm nào trong Qdrant, hỗ trợ \texttt{--chunker legacy} (commit \texttt{c50a651}) hoặc \texttt{--chunker current}, ghi Silver chunks trước rồi mới upsert vector theo đúng bất biến của \texttt{QueueConsumer}.
  \item \path{scripts/backfill_qdrant_year.py}: Cập nhật hàng loạt trường \texttt{year} (lô \texttt{200} phép toán \texttt{set\_payload} mỗi HTTP request) cho các điểm cũ trong Qdrant từ \texttt{published\_date} mà không phải nhúng lại vector.
\end{itemize}

\section{Thuật Toán Phân Mảnh Ngữ Cảnh Đa Tầng (\texttt{AcademicChunker})}
\label{sec:ch2_chunker}

\subsection{Dẫn Xuất Toán Học Cửa Sổ Cân Bằng và Bám Ranh Giới Câu}
\label{sec:ch2_chunker_math}

Ở các phiên bản tiền nhiệm (commit \texttt{c50a651}), chiến lược cắt đoạn dùng cửa sổ trượt kích thước tối đa cố định $S$ và độ gối đầu cố định \texttt{50} từ. Nhược điểm toán học của cách tiếp cận này là khi một mục có độ dài vượt nhẹ bội số hiệu dụng --- chẳng hạn $L = 460$ từ với $S = 450$ --- cửa sổ cuối cùng chỉ còn lại vài chục từ ngắn ngủi ($p_{10} = 103$ từ), nghèo ngữ nghĩa khi vector hóa nhưng vẫn chiếm vị trí ứng viên Top-$k$.

Module \texttt{AcademicChunker} (\path{src/indexing/chunker.py}) giải quyết triệt để vấn đề này bằng thuật toán \texttt{split\_words(words)} với tham số mặc định \texttt{max\_chunk\_words = 300} (được \texttt{QueueConsumer} cấu hình thành $S = 450$ để bao trọn các đoạn giàu công thức toán học), \texttt{overlap\_ratio} $r = 0.15$, giới hạn gối đầu $[o_{\min}, o_{\max}] = [20, 80]$ từ, và ngưỡng gộp đuôi \texttt{min\_tail\_ratio} $\tau = 0.4$:

\begin{enumerate}[label=\textbf{\arabic*.}]
  \item \textbf{Xác định số cửa sổ tối thiểu $n$ và cân bằng độ dài cửa sổ}: Nếu tổng số từ của mục thỏa $L \le S$, giữ nguyên một phân đoạn duy nhất \texttt{[words]}. Ngược lại ($L > S$), số cửa sổ tối thiểu $n \ge 2$ được khởi tạo và lặp tăng dần cho đến khi độ dài thực tế $c \le S$:
\end{enumerate}
\begin{equation}
  n_0 = \max\!\left(2,\; \left\lceil \frac{L - r \cdot S}{S - r \cdot S} \right\rceil\right)
  \label{eq:ch2_chunk_n0}
\end{equation}
\noindent Trong đó: $L$ là tổng số từ của phân mục hiện tại, $S$ là kích thước cửa sổ tối đa (\texttt{max\_chunk\_words}, cấu hình \texttt{450} trong \texttt{QueueConsumer}), $r = 0.15$ là tỷ lệ gối đầu mục tiêu, và $n_0$ là cận dưới số phân đoạn ban đầu.

Tại mỗi vòng kiểm tra số cửa sổ $n \ge n_0$, độ dài ước lượng cân bằng $c_0$, số từ gối đầu động $o$ (tỷ lệ \texttt{15\%} độ dài cửa sổ thực tế thay vì cố định \texttt{50} từ), và độ dài cửa sổ phân bố đều $c$ được tính như sau:
\begin{align}
  c_0 &= \frac{L}{n - r(n - 1)}, \qquad
  o = \min\!\Bigl(o_{\max},\; \max\bigl(o_{\min},\; \operatorname{round}(r \cdot c_0)\bigr)\Bigr), \notag \\
  c &= \left\lceil \frac{L + (n - 1) \cdot o}{n} \right\rceil
  \label{eq:ch2_chunk_balance}
\end{align}
\noindent Trong đó: $n$ là số cửa sổ thử hiện tại, $[o_{\min}, o_{\max}] = [20, 80]$ là biên độ từ gối đầu cho phép, $o$ là số từ gối đầu thực tế sau khi kẹp biên, và $c$ là độ dài mục tiêu mỗi cửa sổ; nếu $c > S$, thuật toán tăng $n \leftarrow n + 1$ và tính lại cho tới khi $c \le S$, bảo đảm kích thước các đoạn đồng đều thay vì dồn phần dư về đoạn cuối.

\begin{enumerate}[label=\textbf{\arabic*.}, start=2]
  \item \textbf{Tính bước trượt đều và Bám ranh giới câu có nhận thức viết tắt học thuật}: Khoảng dịch lý tưởng giữa các cửa sổ và bán kính dung sai tìm kiếm ranh giới câu được xác định bởi:
\end{enumerate}
\begin{equation}
  \text{stride} = \frac{L - c}{n - 1}, \qquad \text{tol} = \max\!\left(5,\; \left\lfloor \frac{o}{2} \right\rfloor\right)
  \label{eq:ch2_chunk_stride}
\end{equation}
\noindent Trong đó: $\text{stride}$ là bước dịch lý tưởng (số từ thực) giữa điểm bắt đầu hai cửa sổ liên tiếp và $\text{tol}$ là bán kính dung sai (tối thiểu \texttt{5} từ) dùng để neo điểm cắt về ranh giới câu tự nhiên.

Hàm \texttt{\_sentence\_starts(words)} quét trước các chỉ mục từ bắt đầu câu mới (từ liền trước kết thúc bằng dấu \texttt{.}, \texttt{!}, \texttt{?} theo sau bởi dấu ngoặc kép/đơn tùy chọn \texttt{\_SENT\_END}), đồng thời \textbf{loại trừ các chữ viết tắt học thuật phổ biến} trong tập \texttt{\_NOT\_SENT\_END = \{"e.g.", "i.e.", "al.", "etc.", "vs.", "cf.", "fig.", "figs.", "eq.", "eqs.", "sec.", "no.", "dr.", "mr.", "ms."\}} cũng như ký hiệu viết tắt một chữ cái hoặc số mục (\texttt{[A-Z]\textbackslash.|\textbackslash d\{1,2\}\textbackslash.}). Với mỗi cửa sổ $k \in \{0, \dots, n-1\}$, điểm bắt đầu lý tưởng $s_{\text{ideal}} = \operatorname{round}(k \cdot \text{stride})$ và kết thúc $e_{\text{ideal}} = s_{\text{ideal}} + c$ được kéo về ranh giới câu gần nhất trong khoảng dung sai \texttt{[lo, hi]} qua hàm \texttt{\_snap()}, đồng thời ràng buộc $e_{\text{prev}} - s \ge o_{\min}$ để không bao giờ mất vùng gối đầu tối thiểu.

\begin{enumerate}[label=\textbf{\arabic*.}, start=3]
  \item \textbf{Gộp đuôi ngắn (\textit{Short-Tail Merge})}: Sau khi bám ranh giới câu, nếu khoảng $(l_s, l_e)$ của cửa sổ cuối cùng ngắn hơn $40\%$ ngưỡng trần và khi nhập vào cửa sổ liền trước $(p_s, p_e)$ vẫn không vượt quá kích thước tối đa $S$:
\end{enumerate}
\begin{equation}
  \bigl(l_e - l_s < 0.4 \cdot S\bigr) \;\land\; \bigl(l_e - p_s \le S\bigr) \;\implies\; [(p_s, p_e), (l_s, l_e)] \leftarrow [(p_s, l_e)]
  \label{eq:ch2_chunk_tail_merge}
\end{equation}
\noindent Trong đó: $(p_s, p_e)$ và $(l_s, l_e)$ lần lượt là cặp chỉ số từ bắt đầu--kết thúc của cửa sổ kế cuối và cửa sổ cuối cùng, còn $0.4$ (\texttt{min\_tail\_ratio}) là ngưỡng tỷ lệ kích hoạt gộp nhằm triệt tiêu triệt để đoạn đuôi vụn.

\subsection{Phân Cấp Ngữ Nghĩa Bốn Tầng và Làm Giàu Tiền Tố}
\label{sec:ch2_chunker_hierarchy}

Áp dụng thuật toán \path{split_words()} trên nền cấu trúc tài liệu, hàm \path{AcademicChunker.chunk_paper()} phân rã mỗi bài báo thành cây phân đoạn 4 tầng:
\begin{itemize}
  \item \textbf{Tầng 1 (Abstract độc lập)}: Tóm tắt của bài báo luôn được đóng gói nguyên vẹn thành chunk \path{{pid}_abs} để đại diện ngữ nghĩa ở mức bài báo.
  \item \textbf{Tầng 2 (Phân mục cấu trúc \path{sections_json})}: Khi bài có mảng section hợp lệ ($\ge 25$ từ/mục), các mục ngắn ($\le S$ từ) giữ nguyên dạng \path{{pid}_sec{idx}}; các mục dài hơn đi qua \path{split_words()} tạo thành \path{{pid}_sec{idx}_p{part}}.
  \item \textbf{Tầng 3 (Regex Heading Fallback / Body Sliding)}: Khi khuyết \path{sections_json} nhưng phần thân có ít nhất \texttt{40} từ, tách theo mẫu tiêu đề Markdown/số mục (Introduction, Methodology, Experiments\ldots) hoặc cắt trượt thành \path{{pid}_body_p{part}}.
  \item \textbf{Tầng 4 (Tách đôi ngữ nghĩa Abstract)}: Nếu bài báo chỉ có đúng 1 chunk tóm tắt với ít nhất \texttt{160} từ, hệ thống tách thêm hai lát cắt \path{{pid}_abs_p1} (\textit{Problem \& Context}) và \path{{pid}_abs_p2} (\textit{Methodology \& Results}) gối đầu $\pm 20$ từ.
\end{itemize}

Trước khi vector hóa, mỗi chunk được gắn tiền tố ngữ cảnh vĩ mô: \texttt{context\_text = f'Paper: "\{title\}" | Section: \{section\_title\}\textbackslash n\textbackslash n\{chunk\_text\}'}. Chuỗi \texttt{context\_text} này được đưa vào \texttt{NomicEmbedder} để vector hóa, giúp embedding của một đoạn phương pháp hoặc thực nghiệm luôn neo chặt vào tên công trình gốc và vị trí phân mục, trong khi \texttt{ChunkStore} lưu riêng trường \texttt{text} sạch sẽ để trình bày trên trích dẫn RAG và tính điểm từ vựng.

\subsection{Kiểm Chứng Thực Nghiệm Phân Đoạn (\texttt{Chunking A/B})}
\label{sec:ch2_chunking_ab}

Để kiểm chứng định lượng hiệu quả của thuật toán phân mảnh mới so với phiên bản cố định cũ (\texttt{c50a651}), kịch bản \path{benchmarks/chunking_ab.py} tiến hành đánh giá đối chứng mù trên cùng một tập \texttt{1,000} bài báo (\texttt{MAX\_WORDS=450}, so sánh ba biến thể \texttt{A\_old\_fixed50}, \texttt{B\_new\_pct15} và \texttt{C\_new\_pct0}). Kết quả chi tiết được tổng hợp tại Bảng~\ref{tab:chunking_ab} ở Chương~\ref{chap:evaluation} cho thấy cấu hình \textbf{\texttt{B\_new\_pct15}} (mặc định hiện hành của \texttt{AcademicChunker}) nâng phân vị $p_{10}$ của độ dài chunk từ \texttt{103} lên \texttt{166} từ ($+61.2\%$, triệt tiêu nhóm đuôi quá ngắn), đồng thời đạt \textbf{Recall@5 = 0.7679} (tăng $+0.0536$ so với \texttt{0.7143} của bản cũ) và \textbf{Chunk Hit@5 = 0.7321}, trực tiếp nâng ngưỡng bao phủ ngữ cảnh khi backend trích chọn Top-5 đoạn trích đưa vào LLM.

\section{Vector Hóa Ngữ Nghĩa, Chỉ Mục Qdrant HNSW và Biến Đổi Đặc Trưng}
\label{sec:ch2_embed_index_fe}

\subsection{Kiến Trúc \texttt{NomicEmbedder} v1.5 (\texttt{MAX\_TOKENS=2048}, Mean Pooling và \texorpdfstring{$L_2$}{L2})}
\label{sec:ch2_nomic}

Mô-đun \path{src/indexing/embedder.py} triển khai bộ sinh vector ngữ nghĩa cục bộ dựa trên mô hình mã nguồn mở \path{nomic-ai/nomic-embed-text-v1.5} (768 chiều)~\cite{nussbaum2024nomic,nomicmodelcard2024} với bốn thiết kế kỹ thuật trọng tâm:
\begin{enumerate}[label=\textbf{\arabic*.}]
  \item \textbf{Cửa sổ Ngữ cảnh Siêu rộng \texttt{MAX\_TOKENS = 2048} (RoPE)}: Khác với các bộ mã hóa BERT truyền thống bị giới hạn cứng ở \texttt{512} tokens --- vốn làm cắt cụt phần đuôi của các chunk \texttt{450} từ chứa khoảng \texttt{650} token văn bản hoặc lên tới \texttt{1,600} token khi dày đặc ký hiệu LaTeX --- kiến trúc \textbf{Rotary Position Embeddings (RoPE)} của Nomic v1.5 (hỗ trợ tối đa \texttt{8,192} tokens) cho phép cấu hình \texttt{MAX\_TOKENS = 2048} bao phủ trọn vẹn 100\% nội dung phân đoạn học thuật.
  \item \textbf{Tối ưu Phần cứng Tự động và Bản vá FP16 (\texttt{\_patch\_extended\_mask})}: Hàm khởi tạo tự động dò thiết bị theo thứ tự ưu tiên Apple Silicon \texttt{mps} $\rightarrow$ NVIDIA \texttt{cuda} (bật bán chính xác \texttt{torch.float16} để tăng gấp đôi thông lượng) $\rightarrow$ \texttt{cpu}. Bản vá \texttt{PreTrainedModel.get\_extended\_attention\_mask} bảo đảm mặt nạ mở rộng kế thừa đúng \texttt{dtype} (\texttt{float16} trên CUDA) trên \texttt{transformers} phiên bản mới, kèm nhánh bắt ngoại lệ tràn VRAM tự động hạ tải sang CPU cho micro-batch lỗi.
  \item \textbf{Tiền tố Bất đối xứng (\textit{Asymmetric Task Prefixing})}: Tuân thủ đặc tả huấn luyện tương phản của Nomic v1.5, văn bản lưu trữ luôn được gắn tiền tố \texttt{"search\_document: "} trước khi làm sạch surrogate, còn truy vấn người dùng được gắn tiền tố \texttt{"search\_query: "}.
  \item \textbf{Mean Pooling có Mặt nạ và Chuẩn hóa Mặt cầu Đơn vị $L_2$}: Từ các trạng thái ẩn lớp cuối $\mathbf{h}_i \in \mathbb{R}^{768}$ và mặt nạ chú ý $m_i \in \{0, 1\}$ tại vị trí token $i \in \{1, \dots, T\}$, vector thô $\mathbf{v}_{\text{raw}}$ và vector chuẩn hóa $\hat{\mathbf{v}}$ được tính theo:
\end{enumerate}
\begin{equation}
  \mathbf{v}_{\text{raw}} = \frac{\sum_{i=1}^{T} m_i \, \mathbf{h}_i}{\max\!\left(10^{-9},\; \sum_{i=1}^{T} m_i\right)}, \qquad \hat{\mathbf{v}} = \frac{\mathbf{v}_{\text{raw}}}{\|\mathbf{v}_{\text{raw}}\|_2}
  \label{eq:ch2_mean_pool}
\end{equation}
\noindent Trong đó: $T \le 2048$ là độ dài chuỗi token sau đệm, $\mathbf{h}_i$ là vector ẩn tại vị trí $i$, $m_i$ là giá trị attention mask ($1$ cho token thật, $0$ cho padding), và $\hat{\mathbf{v}} \in \mathbb{R}^{768}$ là vector sau chuẩn hóa $L_2$ ($\|\hat{\mathbf{v}}\|_2 = 1$). Trên mặt cầu đơn vị này, độ tương đồng Cosine thu gọn thành tích vô hướng tốc độ cao:
\begin{equation}
  \operatorname{CosineSim}(\hat{\mathbf{u}}, \hat{\mathbf{v}}) = \frac{\hat{\mathbf{u}} \cdot \hat{\mathbf{v}}}{\|\hat{\mathbf{u}}\|_2 \|\hat{\mathbf{v}}\|_2} = \hat{\mathbf{u}} \cdot \hat{\mathbf{v}}
  \label{eq:ch2_cosine}
\end{equation}
\noindent Trong đó: $\hat{\mathbf{u}}$ và $\hat{\mathbf{v}}$ lần lượt là hai vector câu hỏi và đoạn trích đã được chuẩn hóa $L_2$ về đơn vị độ dài.

\subsection{Quản Trị Chỉ Mục Qdrant HNSW (\texttt{QdrantManager}) và Connection Pool}
\label{sec:ch2_qdrant}

Lớp \texttt{QdrantManager} (\path{src/indexing/qdrant_manager.py}) đóng vai trò cổng kết nối duy nhất tới collection \path{scientific_papers_gold} trên Qdrant~\cite{qdrant2023vector}:
\begin{itemize}
  \item \textbf{Định danh điểm tất định UUIDv5}: Mỗi \texttt{chunk\_id} được ánh xạ thành mã UUID chuẩn RFC 4122 v5 thông qua namespace DNS để bảo đảm phép ghi \texttt{upsert} (\texttt{PUT /collections/\{col\}/points?wait=true}, lô \texttt{100} điểm) luôn lũy đẳng, không sinh điểm trùng khi chạy lại:
\end{itemize}
\begin{equation}
  \text{point\_id} = \operatorname{UUIDv5}\bigl(\texttt{NAMESPACE\_DNS},\; \text{chunk\_id}\bigr)
  \label{eq:ch2_uuidv5}
\end{equation}
\noindent Trong đó: \texttt{NAMESPACE\_DNS} là hằng số chuẩn UUID (\texttt{uuid.NAMESPACE\_DNS}) khớp đồng nhất giữa Python và dịch vụ TypeScript, còn \texttt{chunk\_id} là định danh duy nhất của phân đoạn trong \texttt{ChunkStore}.

\begin{itemize}
  \item \textbf{Đồ thị HNSW và Chỉ mục Payload Lọc nhanh}: Khi khởi tạo collection mới (\texttt{ensure\_collection}), hệ thống cấu hình kích thước \texttt{768}, khoảng cách \texttt{Cosine} trên đồ thị HNSW~\cite{malkov2018efficient}, và tạo 4 chỉ mục payload phục vụ lọc kết hợp trước khi duyệt đồ thị: \path{primary_category} (\texttt{keyword}), \path{paper_id} (\texttt{keyword}), \path{year} (\texttt{integer}), và \path{section_type} (\texttt{keyword}). Ngoài ra, \texttt{QdrantManager} duy trì kết nối bền vững \texttt{requests.Session()} loại bỏ chi phí bắt tay TLS lặp lại, giúp giảm độ trễ truy vấn Top-5 qua WAN từ hơn \texttt{1,000 ms} xuống \texttt{269.46 ms} (xem Bảng~\ref{tab:ch2_embed_qdrant_params} và các siêu tham số toàn hệ thống tại Bảng~\ref{tab:hyperparameters}).
\end{itemize}

\begin{table}[htbp]
  \centering
  \small
  \caption[Thông số kỹ thuật của NomicEmbedder và QdrantManager]{Thông số kỹ thuật cấu hình của \texttt{NomicEmbedder} và \texttt{QdrantManager}}
  \label{tab:ch2_embed_qdrant_params}
  \setlength{\tabcolsep}{3.5pt}
  \begin{tabularx}{\linewidth}{|>{\raggedright\arraybackslash}p{22mm}|>{\raggedright\arraybackslash}p{36mm}|>{\raggedright\arraybackslash}p{38mm}|L|}
    \hline
    \textbf{Thành phần} & \textbf{Tham số kỹ thuật} & \textbf{Giá trị cấu hình} & \textbf{Ý nghĩa \& Mục đích thiết kế} \\
    \hline
    \textbf{Embedder} & Mô hình nền tảng & \path{nomic-ai/nomic-embed-text-v1.5} & Kiến trúc RoPE hỗ trợ ngữ cảnh dài \texttt{8,192} tokens \\
    \hline
    \textbf{Embedder} & \path{MAX_TOKENS} / Số chiều & \texttt{2,048} tokens / \texttt{768-D} & Bao phủ trọn vẹn chunk \texttt{450} từ kèm công thức toán \\
    \hline
    \textbf{Embedder} & Tiền tố tác vụ & \path{search_document:} / \path{search_query:} & Tối ưu độ tương đồng bất đối xứng tài liệu--câu hỏi \\
    \hline
    \textbf{Embedder} & Chính xác phần cứng & \texttt{MPS} / \texttt{CUDA FP16} / \texttt{CPU} & Batch mặc định \texttt{8} chunks (\texttt{73.0 chunks/s} trên RTX 4060) \\
    \hline
    \textbf{Qdrant} & Collection \& Metric & \path{scientific_papers_gold} / \texttt{Cosine} & Đồ thị HNSW với vector đã chuẩn hóa $L_2$ ($\|\hat{\mathbf{v}}\|_2 = 1$) \\
    \hline
    \textbf{Qdrant} & \path{QDRANT_PAYLOAD_FIELDS} & \texttt{5} trường tối thiểu & \path{chunk_id}, \path{paper_id}, \path{year}, \path{primary_category}, \path{section_type} \\
    \hline
    \textbf{Qdrant} & Payload Indexes & \texttt{keyword} (3), \texttt{integer} (1) & Lọc tức thời theo ngành, mã bài, loại mục và khoảng năm \\
    \hline
  \end{tabularx}
  \ReportNote{Nguồn: Kiểm tra trực tiếp hằng số trong \path{src/indexing/embedder.py}, \path{src/indexing/qdrant_manager.py}, \path{src/storage/chunk_store.py} và benchmark phần cứng ngày 19/10/2026 (\texttt{SNAP-04}).}
\end{table}

\subsection{Kỹ Thuật Biến Đổi Đặc Trưng Nối Sang EDA và Khai Phá Dữ Liệu}
\label{sec:ch2_feature_eng}

Dữ liệu sau khi hoàn tất chu trình thu thập, tích hợp, làm sạch, phân mảnh và đánh chỉ mục ở Chương này được biến đổi trực tiếp thành bốn tập đặc trưng đầu vào cho bước Khám phá Dữ liệu (Chương~\ref{chap:data_understanding}) và 4 Trụ cột Khai phá Dữ liệu cùng Hybrid RAG (Chương~\ref{chap:modeling}), như tổng hợp trong Bảng~\ref{tab:ch2_feature_bridge}:
\begin{enumerate}[label=\textbf{\arabic*.}]
  \item \textbf{Biến đổi Lệch phải Bằng Hàm Logarit $\log(1+\mathbf{x})$} (Pillar 4 --- \path{TrendAnomalyMiner}): Do phân phối số từ (\path{total_words}) và số công thức (\path{total_math_count}) có đuôi phải cực dài, vector đặc trưng 5 chiều $\mathbf{x} = [\text{math}, \text{words}, \text{sections}, \text{authors}, \text{categories}]^\top$ được nén khoảng giá trị trước khi đưa vào \path{IsolationForest}~\cite{liu2008isolation}:
\end{enumerate}
\begin{equation}
  \mathbf{x}_{\log} = \log(1 + \mathbf{x})
  \label{eq:ch2_log1p}
\end{equation}
\noindent Trong đó: $\mathbf{x} \in \mathbb{R}_{\ge 0}^{5}$ là vector đếm cấu trúc nguyên thủy của bài báo (điền khuyết \texttt{0} nếu thiếu), và phép biến đổi $\log(1 + \cdot)$ giữ điểm gốc tại $0$ trong khi nén các giá trị ngoại lai nhiều bậc độ lớn về cùng thang đo ổn định.

\begin{enumerate}[label=\textbf{\arabic*.}, start=2]
  \item \textbf{Mã hóa Giỏ Giao dịch Nhị phân One-Hot (\texttt{TransactionEncoder} --- Pillar 1 \texttt{AssociationRuleMiner})}: Mỗi bài báo $p$ được ánh xạ thành tập giao dịch $T_p$ hợp từ nhãn chuyên ngành (\path{cat:*}) và từ khóa từ điển \texttt{KEYWORD\_VOCAB} (\path{tag:*} trên 25 khái niệm AI/DS); chỉ giữ các giỏ có $|T_p| \ge 2$ rồi mã hóa thành ma trận nhị phân thưa để tính độ hỗ trợ trong thuật toán FP-Growth~\cite{han2000mining}:
\end{enumerate}
\begin{equation}
  \operatorname{Support}(X) = \frac{|\{T_p \in D_{\mathrm{tx}} : X \subseteq T_p\}|}{|D_{\mathrm{tx}}|}
  \label{eq:ch2_support}
\end{equation}
\noindent Trong đó: $D_{\mathrm{tx}}$ là tập hợp các giao dịch hợp lệ ($|T_p| \ge 2$), $I$ là toàn bộ tập mục sau chuẩn hóa, và $X \subseteq I$ là một tập mục phổ biến.

\begin{enumerate}[label=\textbf{\arabic*.}, start=3]
  \item \textbf{Trọng số Đồ thị Hợp tác Phân số (Pillar 3 --- \texttt{ScientificGraphMiner})}: Khi xây dựng mạng lưới đồng tác giả dự phòng (khi thiếu \path{silver/citations.parquet}), để ngăn hiện tượng bùng nổ cạnh từ các bài báo quy mô hàng trăm tác giả, mỗi cặp đồng tác giả $(u, v)$ trong một bài báo có $N_p > 1$ tác giả nhận trọng số phân số tỷ lệ nghịch:
\end{enumerate}
\begin{equation}
  W(u, v) = \sum_{p \in P_{uv}} \frac{1}{N_p - 1}
  \label{eq:ch2_frac_weight}
\end{equation}
\noindent Trong đó: $P_{uv}$ là tập các bài báo có cả hai tác giả $u$ và $v$ cùng đứng tên, và $N_p$ là tổng số tác giả của bài báo $p$.

\begin{enumerate}[label=\textbf{\arabic*.}, start=4]
  \item \textbf{Chiếu Không gian Giảm chiều Centered PCA và Chuẩn hóa Động lượng Thị phần (Pillar 2 \& Pillar 4)}: Vector abstract 768-D chuẩn hóa $L_2$ từ Qdrant được trừ kỳ vọng trung bình, chiếu qua \texttt{PCA(n\_components=2)} và nhân hệ số phóng đại \texttt{5.0} phục vụ fit Gaussian Mixture Model (GMM)~\cite{dempster1977maximum} và lưới mật độ $96 \times 72$; đồng thời số bài theo từng chuyên ngành được chuẩn hóa theo tỷ trọng thị phần từng quý ($\Delta\text{Share}$) để triệt tiêu nhiễu tăng trưởng quy mô chung của kho.
\end{enumerate}

\begin{table}[htbp]
  \centering
  \small
  \caption[Ma trận kết nối đầu ra Chương 2 làm đặc trưng cho Chương 3 và Chương 4]{Ma trận kết nối đầu ra Chương 2 làm đặc trưng đầu vào cho Chương 3 và Chương 4}
  \label{tab:ch2_feature_bridge}
  \setlength{\tabcolsep}{3.5pt}
  \begin{tabularx}{\linewidth}{|>{\raggedright\arraybackslash}p{31mm}|>{\raggedright\arraybackslash}p{46mm}|>{\raggedright\arraybackslash}p{41mm}|L|}
    \hline
    \textbf{Nguồn đầu ra (Chương 2)} & \textbf{Phép biến đổi đặc trưng (\textit{Feature Engineering})} & \textbf{Mô-đun tiêu thụ hạ nguồn} & \textbf{Chương liên kết} \\
    \hline
    \path{silver/papers/} (23 cột ZSTD) & Chiếu cột DuckDB S3 + dẫn xuất \path{math_density} trên 1,000 từ & \texttt{EdaEngine} (3 Lăng kính \& Tương quan) & Chương~\ref{chap:data_understanding} \\
    \hline
    \path{silver/papers/} (\path{categories}, \path{title}, \path{abstract}) & Từ điển 25 mẫu \path{KEYWORD_VOCAB} + One-Hot \texttt{TransactionEncoder} & Pillar 1: \texttt{AssociationRuleMiner} (FP-Growth) & Chương~\ref{chap:modeling} \\
    \hline
    Qdrant \path{section_type="abstract"} (768-D $L_2$) & Trừ trung bình + \path{PCA(2) * 5.0} + \path{normalize()} & Pillar 2: \path{SemanticClusterAnalyzer} (KMeans / GMM / DBSCAN) & Chương~\ref{chap:modeling} \\
    \hline
    \path{silver/papers/} (\path{authors}) & Trọng số phân số $1/(N_p - 1)$ trên cặp \texttt{combinations(authors, 2)} & Pillar 3: \texttt{ScientificGraphMiner} (PageRank / Louvain) & Chương~\ref{chap:modeling} \\
    \hline
    \path{silver/papers/} (5 đặc trưng đếm \& ngày) & Biến đổi $\log(1 + \mathbf{x})$ + chuẩn hóa thị phần $\Delta\text{Share}$ & Pillar 4: \texttt{TrendAnomalyMiner} (Isolation Forest) & Chương~\ref{chap:modeling} \\
    \hline
    \texttt{ChunkStore} (\path{silver/chunks/}) + Qdrant HNSW & Thủy hóa Byte-Range S3 + tiền tố \path{search_query:} & \texttt{RetrievalService} + \texttt{RerankerService} + \texttt{RagService} & Chương~\ref{chap:modeling} \\
    \hline
  \end{tabularx}
  \ReportNote{Nguồn: Tổng hợp thiết kế tích hợp liên chương từ \path{src/mining/} và \path{backend/app/services/} (kiểm chứng ngày 25/11/2026).}
\end{table}
\FloatBarrier
